\subsection{A worked example}\label{sec:example}

Every object of this section can be built by hand on a network
small enough to check line by line. Take one internal metabolite
pool, an import reaction $v_0$ whose upper bound is the parameter
$\theta$ (the external substrate supply), and two pathways out of
the pool: a high-yield pathway $v_1$ limited by an enzyme capacity
$c$, and a low-yield pathway $v_2$ limited by a capacity $d$. The
stoichiometric balance is $v_0 = v_1 + v_2$, and growth is
$\Phi = a\,v_1 + b\,v_2$ with $a > b$: the first pathway converts
each unit of substrate to more biomass than the second. We sweep
$\theta$ upward --- a progressive nutrient repletion --- and watch
the optimal flux map $v^*(\theta)$ and the value function
$\Phi(\theta)$.

Because the high-yield pathway is worth more per unit, the optimum
fills $v_1$ first, and the parameter axis splits into three
chambers:
\begin{center}
\small
\begin{tabular}{lcccl}
\toprule
phase & $v_0$ & $v_1$ & $v_2$ & binding constraints \\
\midrule
I: $\theta \le c$ & $\theta$ & $\theta$ & $0$ & import bound \\
II: $c < \theta \le c + d$ & $\theta$ & $c$ & $\theta - c$ & import, $v_1$ capacity \\
III: $\theta > c + d$ & $c + d$ & $c$ & $d$ & both capacities \\
\bottomrule
\end{tabular}
\end{center}
The fluxes are continuous in $\theta$ everywhere; their slopes are
not. At $\theta = c$ the capacity of the high-yield pathway binds
and the network crosses its first wall: $v_1$ flattens (slope $1
\to 0$) while $v_2$ opens ($0 \to 1$) to carry the overflow. At
$\theta = c + d$ the import bound stops binding --- the network is
saturated downstream --- and $v_0$ flattens. These are exactly the
codimension-one boundary crossings of Definition~\ref{def:mu}: the
atomic measure of $v_1$ carries a point mass $-\delta_{c}$, that of
$v_2$ carries $\delta_c - \delta_{c+d}$, and that of $v_0$ carries
$-\delta_{c+d}$. The curvature is not spread over the sweep; it
sits at the two walls, as point masses with exactly computable
weights.

The value function is $\Phi(\theta) = a\theta$, then $ac + b(\theta
- c)$, then the constant $ac + bd$; its measure is $D^2\Phi = (b -
a)\,\delta_c - b\,\delta_{c+d}$. The coupling identity of
Theorem~\ref{thm:coupling} checks by arithmetic: at the first
wall, $a\,\Delta v_1' + b\,\Delta v_2' = a(-1) + b(+1) = b - a$;
at the second, $a\,(0) + b\,(-1) = -b$, the import reaction having
objective coefficient zero. The objective value sees each wall
exactly as the weighted sum of the flux-layer jumps --- the
contraction made visible, on a network small enough to check by
hand.

The gene metric is equally explicit. If gene $g_1$ catalyzes $v_1$
and gene $g_2$ catalyzes $v_2$, integrating along the sweep gives
$\kmu(g_1) = 1$ and $\kmu(g_2) = 2$ (up to the common $1/T$
normalization of Definition~\ref{def:kmu}). The low-yield pathway
is crossed by both walls --- it opens at the first and closes at
the second --- so it carries twice the rerouting burden of the
high-yield pathway. The bypass, not the main road, is where the
geometry accumulates.

One further substitution sharpens the story. Set $a = b$: the two
pathways now yield identical growth, and the first wall becomes
\emph{objective-invisible}. Every split $v_1 + v_2 = v_0$ with $v_1
\le c$ attains the same $\Phi$, the optimal face is a segment, and
$\Phi(\theta) = a\min(\theta, c + d)$ has no kink at $\theta = c$
at all --- the contraction $a(-1) + a(+1) = 0$ annihilates exactly
this wall, while the flux layer must still choose a vertex. Which
vertex it chooses is the declared tie-break: a stage-3 rule
preferring $v_1$ keeps the kink in $v_1$ at $\theta = c$; a rule
preferring $v_2$ moves the flux-layer kink to $\theta = d$. The
wall structure --- where the degeneracy lives --- is
selection-free; the assignment of the kink to a reaction at a
degenerate face is not. This is, in miniature, why the
lexicographic tie-break is declared as part of the metric
(Remark~\ref{rem:lock}), why the value layer is the canonical
tie-break-free carrier (Proposition~\ref{prop:alex}), and why the
genome-scale rank stability under alternative selections
(\S\ref{sec:v8}) is a measured property of the network rather than
a triviality. Machine verification of every number above --- phase
formulas, slope jumps, the coupling identity, the $\kmu$ values,
and both tie-break variants, by LP re-solves at wall-straddling
parameter values --- is deposited with the artifacts.
